\section{A tree of circle phases}\label{sec:torus-construction}

The second construction again concatenates ordered child words into
rows, but a leaf now records a quantized circle phase. We will choose
the increments between rows so that a half-turn of all leaf phases
separates the resulting words. The choice is random; its output is one
fixed word map with separation at every base phase.

\subsection{Prime groups and row increments}

Fix a sufficiently large integer $k$ and set
\begin{equation}\label{eq:torus-parameters}
 g=16k,\qquad m=gk^2=16k^3,\qquad
 R=k^{8k},\qquad T=k^{52k}.
\end{equation}
Let $\T=\R/\Z$, with $|t|_{\T}$ denoting the circle distance from
$t$ to zero. Use a full ordered $m$-ary tree of depth $k$, and divide
the children of each internal node into $g$ groups of $k^2$ children.
Each leaf phase will be quantized into $R$ intervals, and each internal
word will contain $T$ rows.

Let $\mathcal P_k$ be the primes in $[k^2,k^3]$.
Lemma~\ref{lem:primes} gives $|\mathcal P_k|>100g$ for large $k$.
At each internal node $v$, independently between nodes, choose an
ordered list $r_{v,1},\ldots,r_{v,g}$ of distinct members of
$\mathcal P_k$, uniformly without replacement. Put
\[
 P_v=\prod_{j=1}^g r_{v,j}\leq k^{48k}.
\]
To specify the increment of a child $vi$ in group $j$, choose $z_{vi}$
uniformly from $\{0,\ldots,r_{v,j}-1\}$ and define
\begin{equation}\label{eq:torus-coefficients}
 c_{vi}=1+z_{vi}\frac{P_v}{r_{v,j}}.
\end{equation}
Conditionally on all prime lists, all these residues are independent.
The increment $c_{vi}/P_v$ has the decomposition
\[
 \frac{c_{vi}}{P_v}=\frac1{P_v}+\frac{z_{vi}}{r_{v,j}}.
\]
Thus every child has a common drift of $1/P_v$, together with an
independently chosen step on its group's prime grid. The same formula
isolates the groups arithmetically: if $\ell\ne j$, then
$c_{vi}\equiv1\pmod{r_{v,\ell}}$. At each selected prime, only
coefficients from its own group can differ from one.

Associate a phase $w_u\in\T$ to each leaf $u$, and write
$w\in\T^{m^k}$ for the full vector. For a subtree vector $\omega$,
the notation $\omega+t$ means addition of $t$ to every coordinate of
that vector. At a leaf, put
\[
 X_u(w_u)=\bigl(u,\lfloor Rw_u\rfloor\bigr),
\]
using the representative of $w_u$ in $[0,1)$. At an internal node,
define
\begin{equation}\label{eq:torus-word}
 X_v(\omega)=\prod_{a=0}^{T-1}
 \left(
 X_{v1}\!\left(\omega_{v1}+\frac{a c_{v1}}{P_v}\right)
 \cdots
 X_{vm}\!\left(\omega_{vm}+\frac{a c_{vm}}{P_v}\right)
 \right),
\end{equation}
where $\omega_{vi}$ is the subvector on the leaves below $vi$.
The product concatenates rows in increasing $a$.
At height $h$, the word length is $(mT)^h$. The root map $X=X_o$ has
length
\begin{equation}\label{eq:torus-length}
 N=(mT)^k
\end{equation}
over the alphabet
$\Gamma=\{(u,b):u\text{ a leaf},\ b\in\{0,\ldots,R-1\}\}$,
of size $m^kR$.
For fixed construction data, every position in this word quantizes one
leaf phase plus a fixed offset. Hence the range of $X$ is finite and
$X$ is measurable.

\subsection{Good shifts give separation at every phase}

Consider $X_v(\omega)$ and $X_v(\omega+t)$, where the scalar $t$ is
added to every leaf phase below $v$. In a comparison of source row
$a$ with target row $a'$, child $vi$ receives the relative shift
\[
 t+(a'-a)\frac{c_{vi}}{P_v}.
\]
The tagged-row lemma will separate the two words if enough children
are separated for every row pair. Since each $c_{vi}$ is an integer,
the displayed shift depends only on $a'-a$ modulo $P_v$. This leads
to a recursive condition on $t$ alone, independent of the base phase.
Set
\[
 \epsilon=R^{-1},\qquad \delta=1/64,\qquad
 \gamma=1/128=\delta/2.
\]
Call a shift $t$ eligible if $|t|_{\T}\geq\epsilon$.

\begin{definition}[Good shift]\label{def:good}
At a leaf, $t\in\T$ is good if $|t|_{\T}\geq\epsilon$.
At an internal node $v$, it is good if, for every integer $a$ modulo
$P_v$, at least $\delta m$ children $vi$ have
$t+a c_{vi}/P_v$ good at $vi$.
\end{definition}

The condition asks for at least $\delta m$ children, interpreted as a
real threshold on their integer count. Its deterministic consequence
is the following separation estimate.

\begin{lemma}[Conditional separation]\label{lem:torus-good-separation}
Fix the construction data. If a shift $t$ is good at a node $v$ of
height $h$, then for every phase vector $\omega$ on its subtree,
\[
 r_{(mT)^h}\bigl(X_v(\omega),X_v(\omega+t)\bigr)
 \geq\gamma^h(mT)^h.
\]
\end{lemma}
\begin{proof}
At a leaf, two points in the same half-open interval of the partition
into $R$ intervals have circle distance strictly less than $1/R$.
Thus a good shift changes the symbol for every phase, including when
its distance is exactly $1/R$.

Suppose $h\geq1$ and the statement holds at the children. Their common
word length is $B=(mT)^{h-1}$. Compare any source row $a$ and target
row $a'$. In child $vi$, the relative shift is
$t+(a'-a)c_{vi}/P_v$. It depends only on $a'-a$ modulo $P_v$, because
$c_{vi}$ is an integer. Goodness supplies at least $\delta m$ children
where the induction hypothesis applies, so the sum of the deficits of
the child comparisons is at least
$d=\delta m\gamma^{h-1}B$.

Leaf labels give pairwise disjoint child alphabets, and the child blocks
have equal length and fixed order. Since
$d\leq mB/64\leq(m-1)B$, Lemma~\ref{lem:tagged} gives
\[
 r_{(mT)^h}\bigl(X_v(\omega),X_v(\omega+t)\bigr)
 \geq\frac{Td}{2}
 =\frac{\delta}{2}\gamma^{h-1}(mT)^h
 =\gamma^h(mT)^h.
\]
This proves the induction pointwise in $\omega$.
\end{proof}

It remains to make the one root shift $t=1/2$ good. We prove a statement
for each fixed eligible shift that can be propagated from a child to
its parent. The all-phase conclusion will then follow from the
deterministic lemma.

\subsection{The probability of a good fixed shift}

\begin{proposition}[Goodness of a fixed shift]\label{prop:good-probability}
For all sufficiently large $k$, every node $v$ and every fixed
$t\in\T$ with $|t|_{\T}\geq\epsilon$ satisfy
\[
 \Prb[t\text{ is good at }v]\geq9/10,
\]
where probability is over the construction in the subtree below $v$.
Consequently, there exists a realization of the entire tree for which
\begin{equation}\label{eq:torus-separation}
 r_N(X(w),X(w+1/2))\geq\gamma^kN
 \qquad\text{for every }w\in\T^{m^k}.
\end{equation}
\end{proposition}
\begin{proof}
The assertion is immediate at a leaf. Assume it holds at the children
of an internal node $v$. Fix an eligible $t$ before choosing any data
at $v$, and suppress $v$ from the prime notation. For a row residue
$a$ modulo $P_v$, let
\[
 S_a=\{j:a\not\equiv0\pmod{r_j}\}.
\]

In group $j$, the relative shift has the form
\[
 t+\frac{a c_{vi}}{P_v}
 =t+\frac a{P_v}+\frac{a z_{vi}}{r_j}\pmod1.
\]
When $j\notin S_a$, the last term vanishes. When $j\in S_a$, it
ranges uniformly over the $r_j$-point grid. We therefore treat small
sets $S_a$ by choosing the prime list, and large sets $S_a$ by using
the independent child residues.

For the first case, we show that the parent list has probability at
least $1-g/|\mathcal P_k|$ of satisfying
\begin{equation}\label{eq:prime-list-event}
 |t+a/P_v|_{\T}\geq\epsilon
 \quad\text{for every }a\text{ with }|S_a|\leq k.
\end{equation}
Because $P_v$ is squarefree, the reduced denominator of such a fraction
$a/P_v$ is a product of at most $k$ selected primes, and is therefore
at most $K_0=k^{3k}$. Distinct rationals on the circle with denominators
at most $K_0$ have circle distance at least $1/K_0^2$. Since
$2\epsilon=2k^{-8k}<1/K_0^2$, at most one of all such rationals lies
at distance strictly less than $\epsilon$ from $-t$.

This candidate rational, if it exists, is determined by $t$ before
the list is sampled. It is nonzero because $|t|_{\T}\geq\epsilon$,
so its reduced denominator has a prime factor. Fix one such factor.
Unless this prime belongs to the sampled list, the candidate cannot
equal any of the fractions under consideration. A prime outside
$\mathcal P_k$ is never selected; a prime inside the pool is selected
with probability exactly $g/|\mathcal P_k|$. This proves the claimed
probability of \eqref{eq:prime-list-event}.

Now fix a parent prime list satisfying \eqref{eq:prime-list-event}.
The remaining random data for different children are independent:
each consists of its residue $z_{vi}$ and its entire child subtree.
Fix one row residue $a$ modulo $P_v$.

If $|S_a|\leq k$, a child in a group $j\notin S_a$ sees the shift
\[
 t+\frac{a c_{vi}}{P_v}
 =t+\frac a{P_v}\pmod1,
\]
since $r_j\mid a$. It is eligible by \eqref{eq:prime-list-event}.
By induction, the children in these groups independently have the
required goodness with probability at least $9/10$. There are at least
$(g-k)k^2$ such children.

If $|S_a|>k$, use children in groups $j\in S_a$, of which there are
at least $k^3=m/16$. Multiplication by $a$ permutes the residues
modulo $r_j$, so each relative shift is uniform on a translate of the
$r_j$-point grid. An arc of length $2\epsilon$ contains a fraction at
most $2\epsilon+1/r_j$ of that grid. The shift is therefore eligible
with probability at least $1-2\epsilon-1/r_j$. Conditional on any
eligible value, the independent child subtree is good with probability
at least $9/10$. Each indicated child succeeds with probability at
least $4/5$ for large $k$, and these success events are independent.

In either case, select $n=m/16$ independent indicated children.
Each has success probability at least $4/5$, and
$\delta m=m/64=n/4$. Hoeffding's inequality for independent bounded
variables \cite[Theorem~2]{Hoeffding1963} gives
\[
 \Prb[\text{fewer than }\delta m\text{ successes}]
 \leq\exp\!\left(-2n\left(\frac45-\frac14\right)^2\right)
 =\exp(-121m/3200).
\]
The variables need not have identical distributions. A union bound
over the $P_v\leq k^{48k}$ possible values of $a$, followed by the
probability of failure of the parent event, gives
\[
 \Prb[t\text{ is not good at }v]
 \leq\frac{g}{|\mathcal P_k|}
 +\exp(48k\log k-121m/3200)<\frac1{10}
\]
for all sufficiently large $k$. Independence between different
residues $a$ is unnecessary. This proves the induction.

Apply the statement to the one fixed root shift $t=1/2$, which is
eligible for large $k$. Choose a realization in which it is good.
Lemma~\ref{lem:torus-good-separation} then gives
\eqref{eq:torus-separation} for every base phase. No simultaneous
random choice for all eligible shifts is asserted.
\end{proof}

Fix this realization for the rest of the torus argument. The prime
lists and coefficients now have fixed values. Only their stated size
and congruence properties will be used below.
